\section{Token Incentive System}

\label{sec:incentives}
% ══════════════════════════════════════════════════════════════════════════════

\subsection{ZooRep Distribution}

Contributors earn ZooRep tokens for verified contributions. The reward
function incentivizes both quantity and quality:

\begin{equation}
  \text{ZooRep}(o) = R_{\text{base}} \cdot q_{obs} \cdot
  (1 + \gamma \cdot \text{DataValue}(o))
\end{equation}

where $R_{\text{base}} = 1$ ZooRep for a standard observation, $q_{obs}$
is the quality score, and DataValue reflects the scientific value of the
observation (higher for undersampled regions, taxa, and times). The
scaling factor $\gamma = 2.0$ ensures that targeted observations in
data-sparse areas earn up to 3x the base reward.

\subsection{Challenge Bonus Structure}

\begin{table}[H]
\centering
\caption{ZooRep bonus multipliers for challenge completion.}
\label{tab:bonuses}
\begin{tabular}{lc}
\toprule
\textbf{Achievement} & \textbf{Multiplier} \\
\midrule
Challenge completion & 1.5x \\
Survey mode observation & 2.0x \\
First observation in grid cell & 3.0x \\
Streak bonus (7 consecutive days) & 1.5x \\
Community validation (correct ID) & 1.2x \\
Rare species documentation & 5.0x \\
Expert-confirmed identification & 2.0x \\
\bottomrule
\end{tabular}
\end{table}

\subsection{Governance Integration}

ZooRep tokens integrate directly with Zoo DAO governance (ZOO-GOV-2026):
contributors with sufficient ZooRep can vote on research funding
proposals, participate in peer review, and join domain working groups.
This creates a direct pathway from biodiversity data collection to
research governance, ensuring that the most active contributors have
voice in how their data is used.

\subsection{Anti-Gaming Measures}

\begin{itemize}
  \item \textbf{Duplicate detection}: Submissions with identical GPS
    coordinates, timestamps, and images within 1 hour are flagged.
  \item \textbf{Quality gating}: ZooRep is only distributed after
    quality scoring; low-confidence observations earn minimal rewards.
  \item \textbf{Photo verification}: AI checks that submitted photos
    contain plausible biological subjects.
  \item \textbf{GPS validation}: Observation locations are checked
    against plausible habitat maps.
  \item \textbf{Reviewer oversight}: Random sampling of observations
    for expert review, with penalty for systematic inaccuracy.
\end{itemize}


% ══════════════════════════════════════════════════════════════════════════════
